% !TeX root = ../main.tex
%% Berliner Hochschule für Technik -- Abschlussarbeit

\chapter*{Einführung}
\addcontentsline{toc}{chapter}{Einführung}
\markboth{Einführung}{Einführung}
\begingroup
\raggedbottom
\setlength{\emergencystretch}{1em}

\section*{Problemstellung und Motivation}

Ein Smartphone kann rauchbezogene Bilder zeigen, eine kurze Benennungsaufgabe anbieten und unmittelbar danach das subjektive Rauchverlangen erfassen. Ob aus diesem technisch einfachen Ablauf eine hilfreiche Unterstützung im Alltag entsteht, ist jedoch eine eigenständige wissenschaftliche Frage. Die vorliegende Masterarbeit verbindet deshalb die Entwicklung einer mobilen Studienanwendung mit der Untersuchung eines verhaltensmedizinischen Ansatzes: Beim Cue-Labeling werden sichtbare Merkmale eines Rauchreizes sprachlich benannt. Unter dem Namen CueLens wird untersucht, wie sich dieser Ansatz technisch umsetzen und unter weniger kontrollierten Bedingungen als im Labor empirisch prüfen lässt.

Die gesundheitliche Bedeutung der Tabakentwöhnung begründet den Anwendungskontext. Für Erwachsene stehen unterschiedliche Formen der Unterstützung zur Verfügung; digitale Angebote können deren zeitliche und räumliche Erreichbarkeit ergänzen. Die WHO-Leitlinie zur Tabakentwöhnung von 2024 nennt unter anderem Smartphone-Apps als mögliche Ergänzung anderer Unterstützungsangebote oder als Instrument zur Selbsthilfe. Für Apps und KI-basierte Interventionen bewertet sie die Evidenzsicherheit allerdings als gering. Die Verfügbarkeit einer digitalen Anwendung ist somit noch kein Beleg für ihren gesundheitlichen Nutzen \cite{who_tobacco_cessation_guideline_2024}. Für einen neuen Ansatz sind technische Funktionsfähigkeit, empirischer Effekt und mögliche Versorgungsrelevanz jeweils gesondert zu untersuchen.

Den verhaltensmedizinischen Ausgangspunkt bildet die Studie von Tabibnia et al. Sie untersucht die Übertragung des Affect Labeling auf rauchbezogene Reize. Beim Cue-Labeling wählen Teilnehmende zu einem Bild ein passendes rauchbezogenes Wort; beim Cue-Matching ordnen sie dem Zielbild ein anderes Bild zu. Benannt werden Merkmale des Reizes, nicht der eigene Gefühlszustand. In der Laboruntersuchung mit 50 täglich rauchenden Erwachsenen war das berichtete Verlangen während Cue-Labeling geringer als während Cue-Matching. Der Gesamteffekt war klein; bei älteren Teilnehmenden zeigte sich ein stärkerer Unterschied. Die Ergebnisse begründen die Prüfung einer mobilen Umsetzung, belegen aber noch keinen Nutzen für einen dauerhaften Rauchstopp \cite{tabibnia_cuelabeling_2026}.

Für die Informatik ist dieser Ansatz interessant, weil die zentrale Interaktion mit wenigen, auf mobilen Geräten etablierten Elementen realisierbar ist. Bilder, Auswahlflächen und eine Skala für Selbstberichte lassen sich in kurze Nutzungssituationen integrieren. Eine spätere Einbeziehung selbst aufgenommener Bilder könnte zudem einen Bezug zu konkreten Alltagssituationen herstellen. Dafür müsste die Anwendung relevante Bildinhalte hinreichend zuverlässig erkennen und geeigneten Begriffen zuordnen. Die vermeintlich einfache Benennungsaufgabe führt damit zu Fragen der Bildklassifikation, der Benutzerführung, der Datenverarbeitung und der überprüfbaren Durchführung eines Studienprotokolls.

\section*{Von der Laboraufgabe zur mobilen Anwendung}

Die Übertragung in den Alltag verändert sowohl die Nutzungssituation als auch die Bedingungen der wissenschaftlichen Beobachtung. Im Labor können Darbietung, Bearbeitungszeit und Untersuchungsumgebung eng vorgegeben werden. Außerhalb des Labors unterscheiden sich verfügbare Zeit, Ablenkung, Geräteeigenschaften und Netzverbindung. Auch der Zeitpunkt der letzten Zigarette und die aktuelle Situation können zwischen Messungen variieren. Eine mobile Umsetzung muss daher zwei Anforderungen miteinander vereinbaren: Sie soll einen hinreichend einheitlichen Aufgabenablauf gewährleisten und gleichzeitig eine Teilnahme ermöglichen, die sich ohne unverhältnismäßigen Aufwand in den Alltag einfügt.

Dabei ist zwischen einer alltagsnahen Durchführung und einer Erfassung natürlich auftretender Rauchreize zu unterscheiden. Die Nutzung einer App außerhalb des Labors macht ihre vorgegebenen Bildaufgaben noch nicht zu einer Beobachtung spontan auftretenden Rauchverlangens. Der Studienprototyp verwendet vorbereitete Reizbilder und zugehörige Antwortoptionen. Damit lässt sich der Vergleich der Aufgabenbedingungen strukturieren. Die Klassifikation eigener Bilder ist ein zusätzlicher technischer Entwicklungsschritt, dessen Leistungsfähigkeit unabhängig von diesem Vergleich beurteilt werden muss. Diese Unterscheidung ist für die Reichweite der Ergebnisse wesentlich: Eine Aussage über den Aufgabenvergleich ist weder automatisch eine Aussage über die Erkennung beliebiger Alltagsbilder noch über die Wirkung einer situationsabhängig ausgelösten Intervention.

Die technische Umsetzung beeinflusst zugleich die Qualität der erhobenen Daten. Ein verlorener Studienfortschritt, eine unklare Zuordnung zur Versuchsbedingung oder eine fehlerhaft wiederholte Übertragung kann die Auswertung beeinträchtigen. Deshalb umfasst die Entwicklung mehr als die Darstellung der Aufgaben. Erforderlich sind eine nachvollziehbare Zustandsverwaltung, die dauerhafte Speicherung relevanter Fortschrittsinformationen, definierte Schnittstellen und ein kontrollierter Umgang mit Unterbrechungen. Ebenso muss erkennbar bleiben, welche Messwerte von derselben Person stammen. Erst diese Verbindung von Interaktion, Persistenz und Datenmodell ermöglicht einen personenbezogenen Vergleich der Bedingungen.

Für eine gesundheitsbezogene Studienanwendung gehören Datenschutz und Teilnahmeorganisation zu diesen Kernanforderungen. Die Zusammenführung wiederholter Selbstberichte erfordert eine stabile Kennung, während unmittelbar identifizierende Angaben vor allem für Anmeldung und Abrechnung benötigt werden. Daraus ergibt sich die Aufgabe, beide Zwecke technisch möglichst weit zu trennen. Zugleich müssen Einwilligung, verständliche Information und erreichbarer Support in den Ablauf eingebunden sein. Die Arbeit betrachtet diese Anforderungen als Bestandteile des Systementwurfs. Ein möglichst großer Funktionsumfang ist demgegenüber nachrangig, wenn zusätzliche Funktionen die Nachvollziehbarkeit, Wartbarkeit oder zuverlässige Durchführung gefährden.

Auch die Reizdarbietung selbst verlangt eine vorsichtige Einordnung. Die Anwendung setzt sich gerade mit Bildern auseinander, die für die Teilnehmenden rauchbezogen bedeutsam sein können. Ihre Verwendung lässt sich deshalb nicht allein mit einer erwarteten Entlastung begründen. Die Untersuchung prüft eine Hypothese und stellt kein individuell wirksames Behandlungsangebot in Aussicht. Die Beschreibung der Teilnahme muss diesen experimentellen Charakter erkennen lassen. Einfache Bedienbarkeit und kurze Aufgaben sind dabei praktische Mittel, um die Belastung zu begrenzen, ersetzen aber weder eine verständliche Studieninformation noch die freiwillige Entscheidung zur Teilnahme.

\section*{Zielsetzung und Forschungsfrage}

Ziel der Arbeit ist die Konzeption, Implementierung und empirische Untersuchung einer mobilen Anwendung zur Benennung von Rauchreizen. Der medizinisch-informatische Beitrag besteht darin, einen verhaltensmedizinisch begründeten Aufgabenvergleich in ein praktisch betreibbares Softwaresystem zu übersetzen und die Aussagekraft der damit gewonnenen Daten nachvollziehbar zu beurteilen. Die technische Entwicklung und die empirische Prüfung sind miteinander verbunden, verfolgen jedoch unterschiedliche Erfolgskriterien. Eine zuverlässig arbeitende Anwendung kann wissenschaftlich brauchbare Daten liefern, auch wenn diese keinen Vorteil des untersuchten Ansatzes erkennen lassen.

Die zentrale empirische Forschungsfrage lautet: \emph{Ist das unmittelbar nach der Aufgabenbearbeitung berichtete Rauchverlangen bei Cue-Labeling geringer als bei Cue-Matching, wenn der Aufgabenvergleich außerhalb des ursprünglichen Laborsettings durchgeführt wird?} Die aus der Laborarbeit abgeleitete Erwartung ist ein niedrigerer Wert nach Cue-Labeling. Gegenstand der Untersuchung ist damit ein kurzfristiger Unterschied zwischen zwei Bedingungen. Ohne eine zusätzliche Messung vor der jeweiligen Aufgabe bezeichnet dieser Unterschied keine unmittelbar beobachtete Abnahme gegenüber dem individuellen Ausgangszustand. Ebenso wenig wird damit bereits eine Verringerung des täglichen Zigarettenkonsums oder eine längerfristige Abstinenz erfasst.

Die technische Leitfrage ergänzt diesen Vergleich: \emph{Wie lässt sich ein datensparsames, robustes mobiles System gestalten, das rauchbezogene Zuordnungsaufgaben, wiederholte Selbstberichte und eine perspektivische KI-gestützte Reizerkennung unter den Ressourcenbedingungen realer Android-Endgeräte unterstützt?} Diese Frage betrifft die gesamte Verarbeitungskette. Dazu gehören die Repräsentation von Bildern und Antwortoptionen, die Steuerung des Studienfortschritts, die Übertragung der Berichte und die reproduzierbare Auswertung. Für die KI-Komponente treten die Auswahl geeigneter Modelle, eine konsistente Bildvorverarbeitung und die Bewertung ihrer Vorhersagen und Laufzeiten hinzu.

Für die Bewertung eines solchen Nutzens ist außerdem relevant, ob eine Veränderung für Betroffene wahrnehmbar und in ihrer Situation wertvoll wäre. Diese Unterscheidung wird von Weinfurt für die Interpretation klinischer Forschung hervorgehoben \cite{weinfurt_clinically_meaningful_2019}. Die vorliegende Untersuchung kann hierzu nur einen begrenzten Beitrag leisten. Sie erfasst das subjektive akute Rauchverlangen in einem definierten Aufgabenablauf. Aussagen über einen Einsatz in der Versorgung würden darüber hinaus Untersuchungen zu längerfristigen Endpunkten, unerwünschten Wirkungen und zur Nutzung unter regulären Versorgungsbedingungen voraussetzen.

\section*{Technischer Beitrag und Abgrenzung der KI}

Die technische Arbeit setzt bei der Frage an, welche Informationen aus einem Bild für die Benennungsaufgabe benötigt werden. Der KI-Prototyp behandelt dies als Klassifikationsproblem mit sechs Zielklassen. Neben Aschenbechern, Zigaretten, Zigarettenpackungen, rauchenden Personen und Rauch ist eine Negativklasse vorgesehen. Sie bildet Eingaben ab, die keiner der festgelegten Rauchreizklassen zugeordnet werden sollen. Diese Begrenzung macht die Aufgabe trainierbar und auswertbar, schränkt aber zugleich die Bedeutung einer Vorhersage ein. Das Modell erkennt vordefinierte visuelle Kategorien; es misst weder das Rauchverlangen einer Person noch die individuelle Bedeutung eines abgebildeten Gegenstands.

Untersucht werden eine allgemeine ML-Kit-Baseline sowie angepasste Modelle der Familien MobileNetV3Small und EfficientNet-Lite0. Ein festgelegter Datensplit und eine dokumentierte Vorverarbeitung schaffen die Grundlage für vergleichbare Auswertungen. Die Beurteilung berücksichtigt klassenbezogene Fehlentscheidungen und Kennzahlen wie Präzision, Sensitivität und F1-Score. Ergänzend erfasst der Android-Benchmark die Inferenzlaufzeit. Dadurch wird die Modellwahl als Abwägung zwischen Erkennungsleistung und praktischer Verwendbarkeit auf mobilen Geräten behandelt. Speicherbedarf, Betriebssystemkompatibilität und mögliche Auswirkungen auf den Energieverbrauch bleiben weitere relevante Kriterien; aus einer kurzen Laufzeit allein folgt noch kein gemessener geringer Energiebedarf.

Fehlklassifikationen haben dabei unterschiedliche Folgen für die vorgesehene Interaktion. Wird ein unpassendes Bild als Rauchreiz eingestuft, kann die Anwendung eine sachlich unzutreffende Benennung anbieten. Wird ein relevanter Reiz übersehen, entfällt eine mögliche Benennungsaufgabe. Solche Fehler betreffen sowohl die Bedienbarkeit als auch die inhaltliche Angemessenheit des Ablaufs. Die Negativklasse und die getrennte Betrachtung der einzelnen Reizklassen sind deshalb für die technische Bewertung bedeutsamer als eine alleinige Angabe der Gesamtgenauigkeit. Ein Modell mit günstigen Durchschnittswerten muss nicht für jede Zielklasse gleichermaßen geeignet sein.

Die KI-gestützte Reizerkennung wird von der eigentlichen Studienlogik getrennt betrachtet. Die bereits mit festen Bildern und Antwortoptionen durchführbare Studie ist nicht davon abhängig, dass ein Modell zuvor beliebige Nutzereingaben zuverlässig klassifiziert. Diese Trennung begrenzt die Zahl gleichzeitig veränderlicher Systemkomponenten und erleichtert die Zuordnung beobachteter Probleme. Sie verhindert zugleich, dass die Wirkung einer Benennungsaufgabe mit der Qualität ihrer automatischen Vorbereitung verwechselt wird. Ein späterer integrierter Einsatz müsste beide Komponenten und ihr Zusammenwirken erneut prüfen.

Eine weitere, davon unabhängige Verwendung von KI liegt in der vorgelagerten Erstellung eines Teils der Bildressourcen. Generierte Bilder werden gesichtet und als feste Ressourcen ausgewählt. Diese Verwendung ist von einer während der Nutzung ausgeführten Bildklassifikation zu unterscheiden. Der Ausdruck \enquote{KI-gestützt} bezeichnet in dieser Arbeit somit konkrete technische Teilaufgaben und keine autonome therapeutische Entscheidungsinstanz. Für die Einordnung der Ergebnisse bleibt jeweils maßgeblich, welche Komponente untersucht wurde und welche Daten ihre Bewertung tragen.

Den Schwerpunkt der Anwendungsentwicklung bildet die native Android-App. Ihre Architektur verbindet die Aufgabenoberfläche mit lokaler Fortschrittsspeicherung und einem Backend für Aktivierung, pseudonymisierte Selbstberichte und Studienabschluss. Die im Projekt dokumentierte iOS-Portierung ergänzt diese Umsetzung, verändert aber nicht den Schwerpunkt auf Android-Endgeräten. Für die Beurteilung sind insbesondere wiederaufnehmbare Abläufe, kontrollierte Eingaben, geschützte Übertragung und eine nachvollziehbare Trennung von Identifikations- und Berichtsdaten relevant. Diese Eigenschaften bilden die praktische Voraussetzung dafür, dass der entwickelte Prototyp überhaupt als Instrument einer gesundheitsbezogenen Untersuchung eingesetzt werden kann.

\section*{Empirischer Ansatz und Aussagegrenzen}

Der empirische Vergleich ist als Within-Subject-Design angelegt: Dieselben Personen bearbeiten beide Aufgabenbedingungen. In der App sind zwanzig Studiensituationen mit jeweils fünf Zuordnungsaufgaben und einem anschließenden Selbstbericht vorgesehen. Damit entstehen pro Person zehn Berichte nach Cue-Matching und zehn nach Cue-Labeling. Das Rauchverlangen wird auf einer Skala von 0 bis 100 angegeben. Die mehrtägige App-Studie sieht einen Mindestabstand von drei Stunden zwischen den Situationen vor. Die beabsichtigte Zielgruppe umfasst täglich rauchende Erwachsene im Alter von 30 bis 65 Jahren mit einem Konsum von mindestens zehn Zigaretten pro Tag und den für die Teilnahme erforderlichen sprachlichen und technischen Voraussetzungen.

Die Arbeit dokumentiert außerdem eine eigenständige browserbasierte Umsetzung auf Testable. Auch dort werden beide Bedingungen und insgesamt zwanzig Selbstberichte verwendet. Anders als in der App wechseln die Bedingungen blockweise, und zwischen den Blöcken sind fünfminütige Pausen vorgesehen. Unterschiede in Zeitplan, Sprache, Endgeräten und Präsentation begrenzen die unmittelbare Vergleichbarkeit beider Erhebungswege. Die Testable-Erhebung wird deshalb gesondert ausgewertet. Ihre Ergebnisse sind als Befunde dieses konkreten Ablaufs zu verstehen und nicht ohne Weiteres als Nachweis für die mehrtägige Nutzung der Android-App.

Für die Auswertung werden die Berichte jeder Person zunächst nach Bedingung zusammengefasst. Entscheidend ist anschließend die Differenz zwischen dem individuellen Labeling- und Matching-Mittelwert. Negative Differenzen entsprechen geringerem berichteten Rauchverlangen nach Labeling. Dieses Vorgehen berücksichtigt, dass wiederholte Angaben derselben Person keine unabhängigen Teilnahmen darstellen. Die Zahl der Berichte darf daher nicht mit der Zahl unabhängiger Untersuchungseinheiten verwechselt werden. Neben dem mittleren Unterschied sind seine Unsicherheit und die Streuung zwischen Personen für die Interpretation wesentlich.

Die Aussagekraft hängt darüber hinaus von der Reihenfolge der Aufgaben ab. In der App folgen auf zehn Matching-Durchgänge zehn Labeling-Durchgänge. Damit können zeitliche Veränderungen, Gewöhnung oder andere im Studienverlauf auftretende Einflüsse mit der Bedingung zusammenfallen. Ein Wechsel der Bedingungen wie bei Testable beseitigt seinerseits nicht automatisch alle Reihenfolge- und Übertragungseffekte. Die Untersuchung wird deshalb als explorative Prüfung eingeordnet. Ein beobachteter Unterschied bedarf einer vorsichtigen kausalen Interpretation; ein nicht signifikanter Unterschied belegt umgekehrt weder die Gleichwertigkeit beider Aufgaben noch die generelle Wirkungslosigkeit von Cue-Labeling.

Der wissenschaftliche Ertrag liegt in der Verbindung einer dokumentierten technischen Umsetzung mit einer nachvollziehbaren empirischen Prüfung. Dazu gehören auch Grenzen der Rekrutierung, der Selbstauskunft, des Messinstruments und des jeweiligen Durchführungskontexts. Die Auswertung soll sichtbar machen, welche Schlussfolgerungen die tatsächlich erhobenen Daten erlauben und welche Fragen offenbleiben. So kann die Arbeit sowohl Anhaltspunkte für eine verbesserte Studienplanung als auch konkrete Anforderungen an eine Weiterentwicklung der Anwendung liefern.

\section*{Aufbau der Arbeit}

Kapitel~1 erläutert die verhaltensmedizinischen und technischen Grundlagen sowie verwandte Arbeiten. Es ordnet Cue-Labeling in den Kontext digitaler Interventionen ein und behandelt Messung, Stichprobenplanung und die verwendeten KI-Ansätze. Kapitel~2 beschreibt die rechtlichen Rahmenbedingungen der Teilnahme, des Datenschutzes und der Bildverwendung. Kapitel~3 entwickelt die technische Evaluation der Bildklassifikation; Kapitel~4 betrachtet die Anforderungen an Endgeräte und Betriebssysteme.

Kapitel~5 führt die Implementierung der App und ihrer Web- und Hintergrundkomponenten zusammen. Es beschreibt außerdem die eigenständige Testable-Umsetzung. Kapitel~6 behandelt Planung und Durchführung der App-Studie einschließlich Rekrutierung, Teilnahmeablauf und Support. Kapitel~7 stellt die Auswertung der Selbstberichte und den reproduzierbaren R-Quelltext dar. Kapitel~8 ist dem Fazit und dem Ausblick gewidmet. Der Aufbau folgt damit dem Weg von der fachlichen Begründung über das technische System und die Studiendurchführung bis zur Beurteilung der empirischen Ergebnisse und ihres begrenzten Geltungsbereichs.

\par
\clearpage
\endgroup
